%! Radicals and Rational Exponents — Unit 1, Lesson 1.2 — Student Packet Cover Sheet
\documentclass[10pt]{article}
\usepackage{atda-article}
\usepackage{atda-boxes}
\usepackage{ltablex}
\keepXColumns

\begin{document}

% Full-bleed royal banner
\begin{tikzpicture}[remember picture, overlay]
  \fill[royal] (current page.north west)
    rectangle ([yshift=-0.9in]current page.north east);
\end{tikzpicture}
\vspace{-0.8in}

\begin{center}
  {\color{white}\textbf{\LARGE Algebra, Trigonometry, and Data Analysis}}\\[4pt]
  {\color{mistmid}\large\textbf{Unit 1: Algebraic Expressions and Factoring}}\\[6pt]
  {\color{white}\textbf{\large Lesson 1.2 \quad Radicals and Rational Exponents}}
\end{center}

\vspace{0.22in}
\namedateperiod
\vspace{0.18in}

\begin{learningtargetbox}
\small
\begin{enumerate}[leftmargin=1.5em, itemsep=5pt]
  \item \ldots{} simplify a radical with a \textbf{numeric or algebraic radicand} by pulling out
        every perfect $n$th-power factor.
  \item \ldots{} add, subtract, multiply, and divide radical expressions, including
        \textbf{rationalizing} a denominator, and simplify the result.
  \item \ldots{} convert between a radical and a \textbf{rational exponent} in both directions,
        and evaluate expressions like $27^{2/3}$ without a calculator.
  \item \ldots{} use a square-root model from a real situation to answer a question about that
        situation, then say what the answer means.
\end{enumerate}
\end{learningtargetbox}

\vspace{0.14in}

\begin{tocbox}
\small
\renewcommand{\arraystretch}{1.5}
\begin{tabularx}{\linewidth}{@{} c l X r @{}}
\rowcolor{mist}
\textbf{\#} & \textbf{Component} & \textbf{Description} & \textbf{Score} \\
\midrule
1 & Warm-Up        & Perfect powers, hunting square factors, and one discovery
                                                                    & \blank{1.2cm} \\
2 & Guided Notes   & Simplifying radicals, rational exponents, and the four operations
                                                                    & \blank{1.2cm} \\
3 & Group Activity & The drop tower: a free-fall model built from a square root
                                                                    & \blank{1.2cm} \\
4 & Exit Ticket    & One simplification, one sum, and one conversion to explain
                                                                    & \blank{1.2cm} \\
5 & Homework       & Mixed practice, Kepler's third law, and a look ahead to factoring
                                                                    & \blank{1.2cm} \\
\midrule
  & \multicolumn{2}{r}{\textbf{Total}} & \blank{1.2cm} \\
\end{tabularx}
\end{tocbox}

\vspace{0.14in}

\begin{remindbox}
\small
This lesson covers \texttt{A2.EO.2a--c} --- simplifying radicals, operating on them, and
converting between radicals and rational exponents. Everything here runs on Lesson 1.1's exponent
rules; the only new thing is that the exponents are now fractions. Watch for the two habits that
matter most: \textbf{simplify before you combine}, and \textbf{a root never splits over a sum}.
\end{remindbox}

\end{document}
